\centerline{\bf \Huge Математическая лабиринт}

\begin{flushright}
        \scriptsize{Если хотите добиться невозмодного - учитесь думать.}
\end{flushright}


\section{Исследование 1}

\problem{1}
Рундук, полный рун, весит 32 кролика, а рундук, заполненный рунами наполовину, - 17 кроликов. Сколько весит пустой рундук?
% 2


\problem{2}
На чудо-дереве росли 28 апельсинов и 19 бананов. Каждый день садовник снима- ет с дерева ровно два фрукта. Если снятые фрукты одинаковы, то на дереве появляется новый банан, а если разные — новый апельсин. В конце концов на дереве оказался только один фрукт. Какой именно?
%один банан


\problem{3}
Чтобы открыть сейф, нужно ввести код - число, состоящее из семи цифр: двоек и троек. Сейф откроется, если двоек больше, чем троек, а код делится и на 3, и на 4. Придумайте код, открывающий сейф.

\problem{4}
Решите уравнение $(x : 2 - 3) : 2 - 1 = 3.$

\problem{5}
Какой вопрос нужно задать жителю острова Рыцари и Лжецы, чтобы узнать, живет ли у него дома ручной крокодил?
% Что бы Вы ответили мне минуту назад, если бы я спросил, живет ли у Вас дома ручной зеленый крокодил?

\problem{6}
Даны шесть чисел 21, 22, 23, 25 и 26. Разрешено к любым двум числам прибавить по единице. Можно ли сделать все числа равными?


\newpage

\section{Исследование 2}

\problem{1}
В домике Совуньи имеется 8 розеток, 21 тройник и неограниченный запас утюгов. Какое наибольшее число утюгов Совунья может включить в сеть одновременно?
% 14

\problem{2}
Найдите наибольшее 10-значное число, состоящее из разных цифр, которое де- лится на 4.

% 999999996

\problem{3}
Отрезок, равный 28 см, разделён на три (возможно неравных) отрезка. Расстояние между серединами крайних отрезков равно 16 см. Найдите длину среднего отрезка.

\problem{4}
Существуют ли два таких натуральных числа, что если их сумму умножить на их произ-ведение, то получится 13281329?
% net

\problem{5}
В семье 5 братьев. У каждого из них есть две сестры. Сколько детей в семье?


\problem{6}
На доске были написаны 10 последовательных натуральных чисел. Когда стерли одно из них, то сумма оставшихся оказалась равна 2002. Какое число стёрли?